\chapter{What This Is}

\standfirst{A Smalltalk-80 written in C that runs bytecodes on every core of
the machine at once, and is checked byte-for-byte against Xerox's own 1983
execution traces. This chapter says what that sentence buys you, and what it
costs you.}

\section{The shortest possible description}

Smalltalk-2026 is a complete Smalltalk system: a language, a class library, a
compiler, a virtual machine, and a windowed programming environment in which
the tools are themselves written in the language. If you have used Pharo or
Squeak, almost everything you know transfers. If you have used none of them,
this book assumes nothing.

Two things make it unlike every other Smalltalk you can install.

\subsection{It uses the whole machine}

Every production Smalltalk---Squeak, Pharo, Cuis, VisualWorks, GNU
Smalltalk---schedules its processes \emph{green}, on one operating-system
thread. \ct{Process}, \ct{ProcessorScheduler} and \ct{Semaphore} are all real
and all work, and they buy you concurrency: many logical threads of control,
interleaved. What they cannot buy you is parallelism. Eight cores sit idle
while one runs your image.

Here the object memory and the scheduler are replaced with ones that use all
of them. Smalltalk processes are still green and still cheap, but they are
multiplexed over a pool of native worker threads, and two of them can be
executing bytecodes on two cores in the same nanosecond. Measured on eight
physical cores, with the work fixed and divided:

\begin{center}
\small
\begin{tabular}{@{}lrrr@{}}
\toprule
kernel & 1 worker & 8 workers & speedup\\
\midrule
arithmetic  & 155.1\,ms & 19.9\,ms & 7.79$\times$\\
mandelbrot  & 380.4\,ms & 48.1\,ms & 7.91$\times$\\
intervals   & 124.9\,ms & 27.3\,ms & 4.58$\times$\\
collections & 191.7\,ms & 29.7\,ms & 6.45$\times$\\
\bottomrule
\end{tabular}
\end{center}

You can reproduce that with \ct{make bench}. Part V of this book is about how
to write code that earns those numbers, and---more importantly---about the one
guarantee you lose in exchange for them.

One qualification, stated here rather than left to be discovered: the worker
pool is started by \ct{make bench} and by the parallel test suites, which run
at up to 31 native threads under ThreadSanitizer. It is not yet started by
\ct{./st80 -run image}, so the windowed desktop runs the interpreter on one
thread today. The contract in Chapter~\ref{ch:contract} applies to what you
write regardless, and Chapter~\ref{ch:building} says where the line is.

\subsection{It is checked against 1983}

The Xerox distribution tape carried reference dumps and execution traces
alongside the image itself. This system reproduces them:

\begin{center}
\small
\begin{tabular}{@{}ll@{}}
\toprule
oracle & result\\
\midrule
\ctp{trace2} & 611 of 611 lines, byte for byte\\
\ctp{trace3} & 482-line prefix exact, one documented divergence\\
\ctp{class.oops} & 446 of 446 classes\\
\ctp{method.oops} & 4{,}494 of 4{,}494 methods\\
\bottomrule
\end{tabular}
\end{center}

That matters to you, the person writing Smalltalk on it, for one practical
reason: when this system does something surprising, the surprise is far more
likely to be Smalltalk-80's than a bug. The semantics were not reconstructed
from memory. They were checked against a recording of the real machine.

\section{What you get}

\begin{description}
\item[The 1983 class library, entire.] All 4{,}500-odd methods of the 226
  classes in \ct{sources/} compile and bootstrap into a live image. The
  Browser, Compiler, Inspector and Debugger you use are the library's own.
\item[A language past the Blue Book.] Full closures with non-local return,
  \ct{ensure:} and \ct{ifCurtailed:}, an exception system with \ct{signal},
  \ct{on:do:}, \ct{retry}, \ct{resume:} and \ct{pass}, general pragmas,
  dynamic arrays \ct|{ }|, byte arrays \ct|#[ ]|, block-local temporaries, and
  traits.
\item[Pharo's own code, on Pharo's own tests.] Packages are imported in Tonel
  format and held to their upstream suites: 1{,}177 Pharo tests run unmodified
  against this system and pass.
\item[Concurrency classes 1983 has no equivalent of.] \ct{Mutex},
  \ct{Monitor}, \ct{SharedQueue}, \ct{Promise}, and
  \ct{Processor>>forkParallel:}---and the 1983 library itself audited on
  thirty-one workers, with what it shared put behind a lock or reorganized
  away.
\item[SQL, on every core at once.] PostgreSQL, MySQL, SQLite, Oracle and SQL
  Server through ODBC, with a query builder that finds its own joins, and one
  connection per process rather than N processes taking turns at one socket.
  Chapter~\ref{ch:database}.
\item[JSON, and exactly.] RFC 8259 read and written, with numbers that are
  still exact when they come out the other side---a tenth times ten is one.
  Chapter~\ref{ch:json}.
\item[A screen that is not from 1983.] The display grows to fill the window,
  and the text is Inter---proportional and antialiased---on a system whose
  BitBlt is still one bit per pixel.
\end{description}

\section{What you do not get}

Being straight about this early will save you an afternoon.

\begin{description}
\item[Morphic, Bloc, Spec, Athens.] None of them, and none of them are
  planned. They are single-threaded by construction. The user interface is
  1983's MVC, and Part IV is a guide to using it, because it will not be
  obvious.
\item[Iceberg, Metacello, Monticello.] There is no package manager. Code
  lives in directories of Tonel files, and which ones compose an image is
  stated in a profile. Chapter~\ref{ch:packages} is that story, and it is
  simpler than the one it replaces.
\item[A REPL you can leave running.] There is no \ct{st80} shell that keeps
  an image warm between expressions. You either evaluate one program with
  \ct{-eval}, or you run the windowed image and work inside it.
\item[Same-priority atomicity.] Smalltalk-80 guarantees that a process is
  never preempted by another of the same priority, and forty years of
  Smalltalk code quietly relies on it as a mutual-exclusion primitive. That
  guarantee is deliberately broken here. This is the single most important
  thing in this book, and Chapter~\ref{ch:contract} is about nothing else.
\end{description}

\section{Who this book is for}

It assumes you can use a terminal and have written programs in some language.
It does not assume you have written Smalltalk. Parts II and III teach the
language and library from the beginning; if you already know Smalltalk, read
Chapter~\ref{ch:beyond} for what is different, then skip to Part IV.

\section{How to read it}

\begin{description}
\item[Part I] gets the system built and gets you inside it.
\item[Part II] is the language: syntax, objects, blocks, classes, exceptions.
\item[Part III] is the library: numbers, collections, streams, text, graphics,
  and talking to a database.
\item[Part IV] is the windowed environment, which is where Smalltalk stops
  being a language and starts being a system.
\item[Part V] is concurrency, and is the reason this system exists.
\item[Part VI] is developing on disk: packages, profiles, tests, tools.
\item[Part VII] is one chapter on how the virtual machine works, for when you
  want to know why something is the way it is.
\end{description}

\begin{stnote}[About the examples]
Every expression in this book whose answer is shown was run against an image
built from this source tree, and the answers were copied out rather than typed
in. They are checked in as \ct{manual/examples/ManualExamples.class.st},
and \ct{make verify} in the \ct{manual} directory re-runs them all. Where
this system answers something other than what Pharo would, the book says so
rather than quietly showing Pharo's answer.
\end{stnote}
